%%%%%%%%%%%%%%%%%%%%%%%%%%%%%%%%%%%%%%%%%%%%%%%%%%%%%%%%%%%%%%%%%%%%%%%%%%%%%%%
% Version Control Systems
%
% Author: FOSSEE 
% Copyright (c) 2009, FOSSEE, IIT Bombay
%%%%%%%%%%%%%%%%%%%%%%%%%%%%%%%%%%%%%%%%%%%%%%%%%%%%%%%%%%%%%%%%%%%%%%%%%%%%%%%%

\documentclass[12pt,compress]{beamer}

\mode<presentation>
{
  \usetheme{Warsaw}
  \useoutertheme{infolines}
  \setbeamercovered{transparent}
}

\usepackage[english]{babel}
\usepackage[latin1]{inputenc}
%\usepackage{times}
\usepackage[T1]{fontenc}

% Taken from Fernando's slides.
\usepackage{ae,aecompl}
\usepackage{mathpazo,courier,euler}
\usepackage[scaled=.95]{helvet}

\definecolor{darkgreen}{rgb}{0,0.5,0}

\usepackage{listings}
\lstset{language=bash,
    basicstyle=\ttfamily\bfseries,
    commentstyle=\color{red}\itshape,
  stringstyle=\color{darkgreen},
  showstringspaces=false,
  keywordstyle=\color{blue}\bfseries}

\newcommand{\inctime}[1]{\addtocounter{time}{#1}{\tiny \thetime\ m}}

\newcommand{\typ}[1]{\lstinline{#1}}

\newcommand{\kwrd}[1]{ \texttt{\textbf{\color{blue}{#1}}}  }

\setbeamercolor{emphbar}{bg=blue!20, fg=black}
\newcommand{\emphbar}[1]

%%%%%%%%%%%%%%%%%%%%%%%%%%%%%%%%%%%%%%%%%%%%%%%%%%%%%%%%%%%%%%%%%%%%%%
% DOCUMENT STARTS
\begin{document}

\begin{frame}

\begin{center}
\vspace{12pt}
\textcolor{blue}{\huge Version Control with hg}
\end{center}
\vspace{18pt}
\begin{center}
\vspace{10pt}
\includegraphics[scale=0.95]{../images/fossee-logo.png}\\
\vspace{5pt}
\scriptsize Developed by FOSSEE Team, IIT-Bombay. \\ 
\scriptsize Funded by National Mission on Education through ICT\\
\scriptsize  MHRD,Govt. of India\\
\includegraphics[scale=0.30]{../images/iitb-logo.jpg}\\
\end{center}
\end{frame}

\begin{frame}
  \frametitle{Objectives}
\label{sec-2}
  At the end of this session, you will be able to:
  \begin{itemize}
  \item Clone existing repositories,
  \item Share your repositories with peers,
  \item Use version control for collaborating with your peers,
  \end{itemize}
\end{frame}

\begin{frame}
\frametitle{Pre-requisite}
\textbf{Version Control Using Hg}
	\begin{itemize}
	\item Part 1
	\item Part 2
	\item Part 3
	\end{itemize}
\end{frame}

\section{Collaborating with Mercurial}
\begin{frame}[fragile]
  \frametitle{Cloning Repositories}
  \begin{itemize}
  \item \typ{hg clone SOURCE [DEST]}
  \item All \typ{hg} repositories are self-contained
  \end{itemize}
\end{frame}

\begin{frame}[fragile]
  \frametitle{Sharing Repositories}
  \begin{itemize}
  \item \typ{hg serve}
  \item Can be cloned with \typ{hg clone http://my-ip-address:8000}
  \item We share a central repository; work on our local copies. 
  \item Set write permissions in \typ{.hg/hgrc}
  \end{itemize}
  \begin{lstlisting}
    [web]
    push_ssl=False
    allow_push=*
  \end{lstlisting}
\end{frame}

\begin{frame}
  \frametitle{Sharing Changes}
  \begin{itemize}
  \item Use \typ{hg push} to push your \typ{commits}
    (\typ{changesets}) to the central repository
  \end{itemize}
\end{frame}


\begin{frame}
  \frametitle{Pulling Changes}
  \begin{itemize}
  \item \typ{hg incoming} shows new \typ{changesets} in the server 
  \item To get these \typ{changesets}, we use \typ{hg pull}
  \item These changes do not affect our working directory
  \item \typ{hg parent} shows the parents of the working directory
  \end{itemize}
\end{frame}

\begin{frame}
  \frametitle{Pulling Changes \ldots}
  \begin{itemize}
  \item \typ{hg update} will update the working directory 
    \begin{itemize}
    \item Updates to the \typ{tip} if no revision is specified
    \item \typ{tip} is the most recently added changeset 
    \item Can specify revision number to update to
    \end{itemize}
  \item \typ{hg tip} shows the \typ{tip} of the repository
  \end{itemize}
\end{frame}

\begin{frame}
  \frametitle{Simultaneous Changes}
  \begin{itemize}
  \item The logs of both repositories will be different
  \item The repositories have diverged
  \item \typ{hg push} fails, in such a scenario
  \item \alert{Never, Never, Never, Ever} use \typ{hg push -f}
  \end{itemize}
\end{frame}

\begin{frame}
  \frametitle{Merging}
  \begin{itemize}
  \item Pull and merge, when \typ{abort: push creates new remote
    heads!}
  \item \typ{hg merge} will merge the two diverged heads
  \item \typ{commit} after you have \typ{merged}!
  \end{itemize}
\end{frame}

\begin{frame}
  \frametitle{Simultaneous Changes \ldots}
  \begin{itemize}
  \item \typ{outgoing} shows the \typ{changesets} that will be pushed
  \item \typ{hg push} works!
  \item Look at the `Change graph'!
  \end{itemize}
\end{frame}

\begin{frame}
  \frametitle{Simultaneous Conflicting Changes}
  \begin{itemize}
  \item What if the changes conflict? -- overlapping edits
  \item \typ{hg push} fails; \typ{hg pull}; \typ{hg merge}
  \item You now get a diff view with 3 panes 
    \begin{itemize}
    \item First --- current file
    \item Second --- \typ{changesets} that you pulled
    \item Third --- file before you made your changes
    \end{itemize}
  \item Resolve conflict and save
  \item \typ{hg commit}; \typ{hg push}
  \item Look at the `Change graph'!
  \end{itemize}
\end{frame}

\section{Conclusion}

\begin{frame}
  \frametitle{\alert{Advice}: Work-flow}
  General work-flow
  \begin{itemize}
  \item \typ{pull}; \typ{update}
  \item Make changes
  \item \typ{commit}
  \item If changes on repo, \typ{pull} and \typ{merge}
  \item \typ{push}
  \end{itemize}
  \emphbar{Commit Early, Commit Often}
\end{frame}

\begin{frame}
\frametitle{Summary}
In this tutorial, we have learnt to,
\begin{itemize}
\item Clone repositories, using hg clone,
\item Serve our repositories via http using hg serve,
\item Push changes to a repository using hg push,
\item Use hg incoming
\item Pull changes from a repository using hg pull ,
\item Update the working directory, using hg update,
\item Merge two heads, using hg merge,
\item Resolve conflicts using hg resolve.
\end{itemize}
\end{frame}

\begin{frame}[fragile]
\frametitle{Evaluation}
\begin{enumerate}
\item Mention the easiest way to get started on sharing your repository by
providing a web interface 
\item Suppose Joey and Melissa have made simultaneous changes to the same file
in their own systems. Would the output of hg parents before and after if one of
them pulls in the changes and merges with it?
\item  What are the commands involved in the process of merging changes? 
\end{enumerate}
\end{frame}


\begin{frame}
\frametitle{Solutions}
\begin{enumerate}
\item hg serve
\item No, whenever we've done a merge, hg parents will display two parents
until we hg commit the results of the merge.
\item hg pull, hg merge, hg commit -m "Merged Remote changes"
\end{enumerate}
\end{frame}


\begin{frame}
\frametitle{Additional Reading}
\begin{enumerate}
\item .hgignore
\item hg rollback
\item hg bisect
\item hg backout
\end{enumerate}
\end{frame}

\begin{block}{}
  \begin{center}
  \textcolor{blue}{\Large THANK YOU!} 
  \end{center}
  \end{block}
\begin{block}{}
  \begin{center}
    For more Information, visit our website\\
    \url{http://fossee.in/}
  \end{center}  
  \end{block}
\end{frame}

\end{document}

